% from ega3-1-fr.tex
\oldpage[III]{87}
\begin{proof}
Les relations (\hyperref[III.1.2.3.1-fr]{1.2.3.1}) sont en effet conséquences
immédiates de ce qu'on a vu dans (\hyperref[III.1.2.2-fr]{1.2.2}). D'autre
part, on a
$C^0(\mathfrak U,\mathcal F)=C^1((\mathbf f),M)$~;
$H^0(\mathfrak U,\mathcal F)$ s'identifie par suite au sous-groupe des
$1$-cocycles de $C^1((\mathbf f),M)$~; comme
$M=C^0((\mathbf f),M)$, la suite exacte
(\hyperref[III.1.2.3.2-fr]{1.2.3.2}) n'est autre que celle qui résulte de la
définition des groupes de cohomologie $H^0((\mathbf f),M)$ et
$H^1((\mathbf f),M)$.
\end{proof}

\begin{corollary}[1.2.4]
\label{III.1.2.4-fr}
Supposons que les $X_{f_i}$ soient quasi-compacts et qu'il existe des
$g_i\in\Gamma(U,\mathcal F)$ tels que
$\sum_i g_i(f_i\vert U)=1\vert U$. Alors, pour tout
$(\mathcal O_X\vert U)$-Module quasi-cohérent $\mathcal G$, on a
$H^p(\mathfrak U,\mathcal G)=0$ pour $p>0$~; si en outre $U=X$,
l'homomorphisme canonique (\hyperref[III.1.2.3.2-fr]{1.2.3.2})
$M\to H^0(\mathfrak U,\mathcal F)$ est bijectif.
\end{corollary}

\begin{proof}
Comme par hypothèse les $U_i=X_{f_i}$ sont quasi-compacts, il en est de même
de $U$, et on peut donc se borner au cas où $U=X$~; l'hypothèse entraîne alors
$H^p((\mathbf f),M)=0$ pour tout $p\geq0$
(\hyperref[III.1.1.9-fr]{1.1.9}). Le corollaire résulte alors aussitôt de
(\hyperref[III.1.2.3.1-fr]{1.2.3.1}) et
(\hyperref[III.1.2.3.2-fr]{1.2.3.2}).
\end{proof}

On observera que puisque
$H^0(\mathfrak U,\mathcal F)=H^0(U,\mathcal F)$ (G, II, 5.2.2), on a ainsi
démontré à nouveau (I, 1.3.7) comme cas particulier.

\begin{remark}[1.2.5]
\label{III.1.2.5-fr}
Supposons que $X$ soit un \emph{schéma affine}~; alors les
$U_i=X_{f_i}=D(f_i)$ sont des ouverts affines, ainsi que les
$U_{i_0i_1\ldots i_p}$ (mais $U$ n'est pas nécessairement affine). Dans ce
cas, les foncteurs $\Gamma(X,\mathcal F)$ et
$\Gamma(U_{i_0i_1\ldots i_p},\mathcal F)$ sont exacts en $\mathcal F$
(I, 1.3.11). Si on a une suite exacte
$0\to\mathcal F'\to\mathcal F\to\mathcal F''\to0$ de
$\mathcal O_X$-Modules quasi-cohérents, la suite des complexes
\[
0\longrightarrow C^\bullet(\mathfrak U,\mathcal F')
\longrightarrow C^\bullet(\mathfrak U,\mathcal F)
\longrightarrow C^\bullet(\mathfrak U,\mathcal F'')
\longrightarrow0
\]
est exacte, et donne donc une suite exacte de cohomologie
\[
\cdots\longrightarrow H^p(\mathfrak U,\mathcal F')
\longrightarrow H^p(\mathfrak U,\mathcal F)
\longrightarrow H^p(\mathfrak U,\mathcal F'')
\xrightarrow{\partial}H^{p+1}(\mathfrak U,\mathcal F')
\longrightarrow\cdots.
\]
D'autre part, si on pose $M'=\Gamma(X,\mathcal F')$,
$M''=\Gamma(X,\mathcal F'')$, la suite
$0\to M'\to M\to M''\to0$ est exacte~; comme
$C^\bullet((\mathbf f),M)$ est un foncteur exact en $M$, on a aussi la suite
exacte de cohomologie
\[
\cdots\longrightarrow H^p((\mathbf f),M')
\longrightarrow H^p((\mathbf f),M)
\longrightarrow H^p((\mathbf f),M'')
\xrightarrow{\partial}H^{p+1}((\mathbf f),M')
\longrightarrow\cdots.
\]
Cela étant, comme le diagramme
\[
\xymatrix{
0\ar[r]&C^\bullet(\mathfrak U,\mathcal F')\ar[r]\ar[d]
&C^\bullet(\mathfrak U,\mathcal F)\ar[r]\ar[d]
&C^\bullet(\mathfrak U,\mathcal F'')\ar[r]\ar[d]&0\\
0\ar[r]&C^\bullet((\mathbf f),M')\ar[r]
&C^\bullet((\mathbf f),M)\ar[r]
&C^\bullet((\mathbf f),M'')\ar[r]&0
}
\]
